% Section 5.3, from lectures/L05/README.md: the objectives, the evaluation questions, and the next
% lecture.

\needspace{10\baselineskip}
\section{Review}
\label{c5:sec:review}

\subsection*{What you should be able to do now}
After this chapter, you should:
\begin{itemize}
  \item Have created a concrete subclass that inherits the existing dense layer interface.
  \item Have implemented all methods, including \code{initParams()}, \code{feedforward()},
    \code{backpropagate()}, and \code{optimize()}.
  \item Have created a fully functioning dense layer implementation.
  \item Have tested the implementation with a complete neural network (the one created in the
    exercises in Chapter~\ref{c4:ch:network}).
\end{itemize}

\needspace{6\baselineskip}
\subsection*{Questions to test yourself}
You should be able to explain:
\begin{enumerate}
  \item What's the difference between \code{ml::dense_layer::Interface} and your concrete class,
    and why is this split beneficial?
  \item Which activation functions does your implementation support, and when is each one
    suitable?
  \item Can you connect the mathematical formulas for backpropagation to your implementation; what
    does each computation step in the code represent?
  \item How do you test that the implementation is correct, and how do you know when the network
    is sufficiently trained?
\end{enumerate}

\needspace{6\baselineskip}
\subsection*{Looking ahead}
Chapter~\ref{c6:ch:cnn} begins Part~II of the book, and introduces convolutional neural networks
(CNNs) for image classification.
